\documentclass{article}

\usepackage{fullpage}
\usepackage{url}
\usepackage{bbm}
\usepackage{graphicx}
\usepackage{amsmath}
\usepackage{physics}
\usepackage{mathtools}
\usepackage{bbold}
\usepackage{slashed}
\usepackage{pxfonts}
\usepackage{multicol}
\usepackage{mathptmx}
\usepackage{simplewick}
\usepackage{placeins}

\usepackage{lmodern}
\usepackage[T1]{fontenc}
\usepackage[fontsize=13.5pt]{fontsize}

\DeclareMathOperator{\erfi}{erfi}
\newcommand{\R}{\ensuremath{\mathbb{R}}}
\newcommand{\C}{\ensuremath{\mathbb{C}}}

\title{Quantum Field Theory I Excercise 5}
\author{Idan Stark}
\date{\today}

\begin{document}
\maketitle

\section{The LSZ reduction formula}
\subsection{The third region pole}
Following the same formulation as for the first integral, we start from
\begin{equation}
    III = \int_{-\infty}^{T^{-}}dx^0 \int d^3 xe^{ip^0x^0 - i\boldsymbol{p}\cdot{\boldsymbol{x}}}
    \bra{\Omega}T\left[\phi(x)\phi(x_2)\cdots\phi(x_{n+m})\right]\ket{\Omega}
\end{equation}
The key difference between the $I$ region and the $III$ region is the position of $x^0$ in the time ordering operator. Where in class we used the fact that it comes before all the other fields, this time it will come last.
Plugging in the identity operator:
\begin{equation}
    \boldsymbol{1} = \sum_{\lambda}\int\frac{d^3q}{(2\pi)^3}\frac{1}{2\epsilon_q(\lambda)}\ket{\lambda_q}\bra{\lambda_q}
\end{equation}
We get:
\begin{equation}
    III = \int_{-\infty}^{T^{-}}dx^0 \int d^3 xe^{ip^0x^0 - i\boldsymbol{p}\cdot\boldsymbol{x}}
    \sum_{\lambda}\int\frac{d^3q}{(2\pi)^3}\frac{1}{2\epsilon_q(\lambda)}
    \bra{\Omega}T\left[\phi(x_2)\cdots\phi(x_{n+m})\right]\ket{\lambda_q}\bra{\lambda_q}\phi(x)\ket{\Omega}
\end{equation}
Using a boost Lorentz transformation to move to the rest frame:
\begin{equation}
\begin{gathered}
    \bra{\lambda_q}\phi(x)\ket{\Omega} = e^{i\boldsymbol{q}\cdot\boldsymbol{x} + i\epsilon_q(\lambda)x^0}\bra{\lambda_0}\phi(0)\ket{\Omega}
\end{gathered}
\end{equation}
The spatial integral becomes a delta function $\delta^{(3)}(\boldsymbol{p} - \boldsymbol{q})$. We can then integrate over $q$ using the $\delta$ function to get:
\begin{equation}
\begin{gathered}
    III =  \sum_{\lambda}\int_{-\infty}^{T^{-}}dx^0 e^{i\left(p^0 + \epsilon_p(\lambda)\right)x^0}
    \frac{1}{2\epsilon_p(\lambda)}
    \bra{\Omega}\phi(0)\ket{\lambda_0}\bra{\lambda_p}T\left[\phi(x_2)\cdots\phi(x_{n+m})\right]\ket{\Omega}
\end{gathered}
\end{equation}
This integral diverges, so let us add $i\epsilon$ to the exponent:
\begin{equation}
\begin{gathered}
    \int_{-\infty}^{T^{-}}dx^0 e^{i\left(p^0 + \epsilon_p(\lambda)\right)x^0} \rightarrow \int_{-\infty}^{T^{-}}dx^0 e^{i\left(p^0 + \epsilon_p(\lambda) + i\epsilon\right)x^0} = \frac{-ie^{i\left(p^0 + \epsilon_p(\lambda) + i\epsilon\right)T^{-}}}{p^0 + \epsilon_p(\lambda) + i\epsilon}
\end{gathered}
\end{equation}
And so the total result is:
\begin{equation}
    III = \sum_{\lambda}
    \frac{1}{2\epsilon_p(\lambda)}
    \frac{-ie^{i\left(p^0 + \epsilon_p(\lambda) + i\epsilon\right)T^{-}}}{p^0 + \epsilon_p(\lambda) + i\epsilon}
    \bra{\Omega}\phi(0)\ket{\lambda_0}\bra{\lambda_p}T\left[\phi(x_2)\cdots\phi(x_{n+m})\right]\ket{\Omega}
\end{equation}
This has a pole at $p^0 = -\epsilon_p = -\sqrt{p^2 + m^2}$ as expected.

\newpage
\section{The optical theorem}
\subsection{$2\rightarrow2$ scatteing in $\phi^4$ theory}
Let us start by writing the full integral of states on the RHS. For two particles with initial momenta $p_A, p_B$ and final momenta $p_C, p_D$, we must integrate over $p_C$ and $p_D$, using the reccomended normlization:
\begin{equation}
    \int \frac{d^3\vec{p}_C}{(2\pi)^3} \frac{d^3\vec{p}_D}{(2\pi)^3} \frac{1}{4E_CE_D} \left|\mathcal{M}(\vec{p}_A, \vec{p}_B \rightarrow \vec{p}_C, \vec{p}_D)\right|^2 (2\pi)^4\delta^{(4)}(p_A + p_B - p_C - p_D)
\end{equation}
But, since the particles are identical, This integral actually double-counts each state. To fix that, lets add a factor $\frac{1}{2}$:
\begin{equation}
    \frac{1}{2}\int \frac{d^3\vec{p}_C}{(2\pi)^3} \frac{d^3\vec{p}_D}{(2\pi)^3} \frac{1}{4E_CE_D} \left|\mathcal{M}(\vec{p}_A, \vec{p}_B \rightarrow \vec{p}_C, \vec{p}_D)\right|^2 (2\pi)^4\delta^{(4)}(p_A + p_B - p_C - p_D)
\end{equation}
Since the LHS is taken to second order, the same must apply here. The only way to achieve that is to take the first order in $\mathcal{M}$, giving:
\begin{equation}
    \frac{\lambda_p^2}{2}\int \frac{d^3\vec{p}_C}{(2\pi)^3} \frac{d^3\vec{p}_D}{(2\pi)^3} \frac{1}{4E_CE_D} (2\pi)^4\delta^{(4)}(p_A + p_B - p_C - p_D)
\end{equation}
Now, let us use the spatial $\delta$ function to solve one of the integrals, let's say $\vec{p}_D = \vec{p}_A + \vec{p}_B - \vec{p}_C$.
\begin{equation}
    \frac{\lambda_p^2}{2}\int \frac{d^3\vec{p}_C}{(2\pi)^3} \frac{2\pi \delta(E_A + E_B - E_C - E_D)}{4E_CE_D}
\end{equation}
Now, moving the to center of mass system, each pair of particles, incoming and outgoing, will have opposite momenta and equal energy
\begin{equation}
\begin{gathered}
    \vec{p}_A = -\vec{p}_B \equiv \vec{p}_i
    \quad
    E_A = E_B \equiv E_i
    \\
    \vec{p}_C = -\vec{p}_D \equiv \vec{p}_f
    \quad
    E_C = E_D \equiv E_f
\end{gathered}
\end{equation}
So, the delta function is reduced to $\delta(2E_i - 2E_f)$, which is obviously equivilant to $\delta(p_i - p_f)$ (a 1D $\delta$ over the 3-momentum magnitude) up to a Jacobean $E_f/2p_f$. Since the integral is spherically symmetric:
\begin{equation}
\begin{gathered}
    \frac{\lambda_p^2}{2}\int \frac{d^3p_f}{(2\pi)^3} \frac{2\pi \delta(p_i - p_f)}{8E_fp_f} =
    \frac{\lambda_p^2}{64\pi^2}\int \frac{\delta(p_i - p_f)}{p_f\sqrt{p_f^2 + m^2}}d^3p_f = 
    \frac{\lambda_p^2}{16\pi}\int \frac{\delta(p_i - p_f)}{\sqrt{p_f^2 + m^2}}p_f dp_f = \\ =
    \frac{\lambda_p^2}{16\pi}\sqrt{\frac{p_i^2}{p_i^2 + m^2}}
\end{gathered}
\end{equation}
Now, we can rewrite it based on $s$:
\begin{equation}
    s = (p_A + p_B)^2 = 4E_i^2 = 4\left(p_i^2 + m^2\right)
\end{equation}
So:
\begin{equation}
    p_i^2 = \frac{s}{4} - m^2
\end{equation}
This gives:
\begin{equation}
\begin{gathered}
    \frac{\lambda_p^2}{16\pi}\sqrt{\frac{p_i^2}{p_i^2 + m^2}} =
    \frac{\lambda_p^2}{16\pi}\sqrt{\frac{\frac{s}{4} - m^2}{\frac{s}{4}}} =
    \frac{\lambda_p^2}{16\pi}\sqrt{1 - \frac{4m^2}{s}}
\end{gathered}
\end{equation}
Whihc is exactly twice what we got for $\Im[\mathcal{M}(\vec{p}_A, \vec{p}_B \rightarrow \vec{p}_A, \vec{p}_B)]$, as expected from the optical theorem.
\subsection{$2\rightarrow2$ scatteing in $\phi^3$ theory}
In the case of the $\phi^3$ theory, the RHS will include all transitions from a 2-particle state to some other final state of order $\lambda_p$, which include just the $2 \rightarrow 1$ transition.
\newline
The LHS will get contribution from the s, t, and u channels. Each channel will give a Faynman propegator with the matching Mandelstam variable.
\begin{equation}
    iM = - \frac{\lambda^2}{2}\left[\frac{i}{s - m^2 + i\epsilon} + \frac{i}{t - m^2 + i\epsilon} + \frac{i}{u - m^2 + i\epsilon}\right]
\end{equation}
So, for $\mathcal{M}$ to have an imaginary part, the propegators must have a real part. This will one of the Mandelstam variables equals to the total mass:
\begin{equation}
    s = m^2 \quad \text{or} \quad t = m^2 \quad \text{or} \quad u = m^2
\end{equation}
Considering only the cases where the momenta are different, this is not the case, and therefore the LHS vanishes.
\newline
This mathces the RHS on the case where $m \ne 0$. Then, the $2 \rightarrow 1$ process cannot conserve mass, and therefore is not allowed. This means that for $m \ne 0$, the RHS vanishes.
\newline
For $m = 0$, the LHS then contributes if one of the Mandelstam variables vanishes. For instance, taking the $s$ variable gives:
\begin{equation}
    s = (p_1 + p_2)^2 = p_1^2 + p_2^2 + 2p_1p_2 = 2m^2 + 2p_1p_2 = 2p_1p_2 = 0
\end{equation}
This means the two 4-momenta must be perpendicular in Minkowski space, i.e. that their product is zero. For massless particles, where $E = p$, this happens if they are in the same or opposite directions. In that case, the matching propegator will give $\frac{\lambda^2}{2}$, and the RHS will give the matching $2 \rightarrow 1$ diagram, resulting in $\lambda^2$ as expected.

\newpage
\section{Two-loop renormalization of the $\phi^4$ theory in four dimensions}
\subsection{Dimensional analysis}
The $-i\Sigma(p^2)$ has dimensions $\left[-i\Sigma(p^2)\right] = 2$. If $m_0 = 0$, the only other scale in the propegator of dimension two is $\left[p^2\right] = 2$. This means that up to some dimensionless function $h$ of $p^2\Lambda^2$, where we divide by the regularization scale $\Lambda$ to make $p^2$ unitless, the sum is
\begin{equation}
    -i\Sigma\left(p^2\right) = p^2 h\left(\frac{p^2}{\Lambda^2}\right)
\end{equation}
\subsection{The sunset diagram contribution to the 1PI propegator $\Sigma$}
Using our trusted ol' Feynman rules, we can write the contribution of ths sunset diagram to the 1PI propegator. Let us name the three internal momenta $k_1, k_2,$ and $k_3$, and notice there is a symmetry factor $6$:
\begin{equation}
\begin{gathered}
    \Sigma_\Theta = -\frac{\lambda^2}{6} \int \frac{d^4k_1}{(2\pi)^4} \frac{d^4k_2}{(2\pi)^4} \frac{d^4k_3}{(2\pi)^4} \frac{i}{k_1^2 + i\epsilon} \frac{i}{k_2^2 + i\epsilon} \frac{i}{k_3^2 + i\epsilon} \times \\ \times
    (2\pi)^4\delta^{(4)}\left(p - k_1 - k_2 - k_3\right) = \\ =
     -\frac{\lambda^2}{6} \int \frac{d^4k_1}{(2\pi)^4} \frac{d^4k_2}{(2\pi)^4}\frac{i}{k_1^2 + i\epsilon} \frac{i}{k_2^2 + i\epsilon} \frac{i}{(p - k_1 - k_2)^2 + i\epsilon}
\end{gathered}
\end{equation}
This integral is of order $\left(k^4\right)^2\left(k^{-2}\right)^3 \sim k^2$ and thus diverges quadratically. Moving to a general dimension $d$:
\begin{equation}
\begin{gathered}
    \Sigma_\Theta = -\frac{\lambda^2}{6} \int \frac{d^dk_1}{(2\pi)^d} \frac{d^dk_2}{(2\pi)^d}\frac{i}{k_1^2 + i\epsilon} \frac{i}{k_2^2 + i\epsilon} \frac{i}{(p - k_1 - k_2)^2 + i\epsilon}
\end{gathered}
\end{equation}
Now, let us use Feynman's trick with two Feynman variables:
\begin{equation}
    \frac{1}{ABC} = 2\int_0^1 dx \int_0^{1-x} dy \frac{1}{\left[Ax + By + C(1 - x - y)\right]^3}
\end{equation}
We can write:
\begin{equation}
    A = k_1^2 \qquad B = k_2^2 \qquad C = (p - k_1 - k_2)^2
\end{equation}
To get:
\begin{equation}
\begin{gathered}
    Ax + By + C(1 - x - y) = \\ =
    k_1^2x + k_2^2y + (p - k_1 - k_2)^2(1 - x - y) = \\ =
    k_1^2x + k_2^2y + \left[
        (p - k_2)^2 + k_1^2 + 2(p - k_2)k_2
    \right](1 - x - y) = \\ =
    (1 - y)k_1^2 + 2(1 - x - y)(k_2 - p)k_1 + yk_2^2 + (1 - x - y)(k_2 - p)^2 = \\ =
    (1 - y)\left[
        k_1^2 + 2\frac{1 - x - y}{1 - y}(k_2 - p)k_1 + \frac{y}{1 - y}k_2^2 + \frac{1 - x - y}{1 - y}(k_2 - p)^2
    \right]
\end{gathered}
\end{equation}
So, we can define
\begin{equation}
    \alpha = \frac{1 - x - y}{1 - y} = 1-  \frac{x}{1 - y}
    \qquad
    \beta = \frac{y}{1 - y}
\end{equation}

And then the integral becomes:
\begin{equation}
\begin{gathered}
    \Sigma_\Theta = i\frac{\lambda^2}{3} \int \frac{d^dk_1}{(2\pi)^d} \frac{d^dk_2}{(2\pi)^d} \int_0^1 dx \int_0^{1-x} \frac{dy}{(1 - y)^3} \times \\ \times \frac{1}{\left[
        k_1^2 + 2\alpha(k_2 - p)k_1 + \beta k_2^2 + \alpha(k_2 - p)^2
    \right]^3}
\end{gathered}
\end{equation}
Completing the squre by defining $l_1 = k_1 + \alpha(k_2 - p)$ gives:
\begin{equation}
\begin{gathered}
    \Sigma_\Theta = i\frac{\lambda^2}{3} \int_0^1 dx \int_0^{1-x} \frac{dy}{(1 - y)^3} \int \frac{d^dk_2}{(2\pi)^d} \frac{d^dl_1}{(2\pi)^d} \frac{1}{\left[
        l_1^2 + \beta k_2^2 + (\alpha - \alpha^2)(k_2 - p)^2
    \right]^3} = \\ =
    i\frac{\lambda^2}{3} \frac{\Gamma\left(3 - \frac{d}{2}\right)}{\Gamma(3)(4\pi)^{d/2}} \int_0^1 dx \int_0^{1-x} \frac{dy}{(1 - y)^3} \int \frac{d^dk_2}{(2\pi)^d} \left(
        \frac{1}{\beta k_2^2 + (\alpha - \alpha^2)(k_2 - p)^2}
    \right)^{3 - \frac{d}{2}}
\end{gathered}
\end{equation}
Again, let us complete the square:
\begin{equation}
\begin{gathered}
    \beta k_2^2 + (\alpha - \alpha^2)(k_2 - p)^2 = (\beta +- \alpha - \alpha^2)k_2^2 - 2(\alpha - \alpha^2)pk_2 + (\alpha - \alpha^2)p^2 = \\ =
    (\beta + \alpha - \alpha^2)\left[
        k_2^2 - 2\frac{\alpha - \alpha^2}{\beta + \alpha - \alpha^2}pk_2 + \frac{\alpha - \alpha^2}{\beta + \alpha - \alpha^2}p^2
    \right]
\end{gathered}
\end{equation}
Defining
\begin{equation}
    \gamma = \frac{\alpha - \alpha^2}{\beta + \alpha - \alpha^2}
\end{equation}
We can define $l_2 = k_2 - \gamma p$ and write the integral as:
\begin{equation}
\begin{gathered}
    \Sigma_\Theta = i\frac{\lambda^2}{3} \frac{\Gamma\left(3 - \frac{d}{2}\right)}{\Gamma(3)(4\pi)^{d/2}} \int_0^1 dx \int_0^{1-x} \frac{dy}{(1 - y)^3(\beta + \alpha + \alpha^2)^{3 - \frac{d}{2}}} \times \\ \times \int \frac{d^dk_2}{(2\pi)^d} \left(
        \frac{1}{k_2^2 - 2\gamma pk_2 + \gamma p^2}
    \right)^{3 - \frac{d}{2}} = \\ =
    i\frac{\lambda^2}{3} \frac{\Gamma\left(3 - \frac{d}{2}\right)}{\Gamma(3)(4\pi)^{d/2}} \int_0^1 dx \int_0^{1-x} \frac{dy}{(1 - y)^3(\beta + \alpha + \alpha^2)^{3 - \frac{d}{2}}} \times \\ \times \int \frac{d^dl_2}{(2\pi)^d} \frac{1}{\left[l_2^2 + (\gamma - \gamma^2) p^2\right]^{3 - \frac{d}{2}}} = \\ =
    i\frac{\lambda^2}{3} \frac{\Gamma\left(3 - \frac{d}{2}\right)}{\Gamma(3)(4\pi)^{d/2}}  \frac{\Gamma\left(3 - \frac{d}{2} - \frac{d}{2}\right)}{\Gamma\left(3 - \frac{d}{2}\right)(4\pi)^{d/2}} \int_0^1 dx \int_0^{1-x} \frac{dy}{(1 - y)^3(\beta + \alpha + \alpha^2)^{3 - \frac{d}{2}}} \times \\ \times \left(\frac{1}{(\gamma - \gamma^2) p^2}\right)^{3 - \frac{d}{2} - \frac{d}{2}} = \\ =
    i\frac{\lambda^2}{3} \frac{\Gamma\left(3 - d\right)}{\Gamma(3)(4\pi)^{d}} \int_0^1 dx \int_0^{1-x} \frac{dy}{(1 - y)^3(\beta + \alpha + \alpha^2)^{3 - \frac{d}{2}}} \left((\gamma - \gamma^2) p^2\right)^{d - 3} = \\ =
    i\frac{\lambda^2}{3} \frac{\Gamma\left(3 - d\right)}{\Gamma(3)(4\pi)^{d}} p^{2(d - 3)} \int_0^1 dx \int_0^{1-x} \frac{dy}{(1 - y)^3(\beta + \alpha + \alpha^2)^{3 - \frac{d}{2}}} \left(\gamma - \gamma^2\right)^{d - 3}
\end{gathered}
\end{equation}
Noting that the remining $dx dy$ integral is simply a number, we can denoting it for simplicity
\begin{equation}
    I_d = \int_0^1 dx \int_0^{1-x} \frac{dy}{(1 - y)^3(\beta + \alpha - \alpha^2)^{3 - \frac{d}{2}}} \left(\gamma - \gamma^2\right)^{d - 3}
\end{equation}
And get:
\begin{equation}
    \Sigma_\Theta = i\frac{\lambda^2}{3} \frac{\Gamma\left(3 - d\right)}{\Gamma(3)(4\pi)^{d}} p^{2(d - 3)} I_d
\end{equation}
Now, taking $d = 4 - \epsilon$:
\begin{equation}
    \Sigma_\Theta = i\frac{\lambda^2}{3} \frac{\Gamma\left(-1 + \epsilon\right)}{\Gamma(3)(4\pi)^{4 - \epsilon}} p^{2(1 - \epsilon)} I_{4 - \epsilon} = 
    i\frac{\lambda^2}{3} \frac{\Gamma(\epsilon)}{2(-1 + \epsilon)(4\pi)^{4 - \epsilon}} p^2 p^{-2\epsilon} I_{4 - \epsilon}
\end{equation}
And using
\begin{equation}
    p^{-2\epsilon} = e^{-2\epsilon\ln(p)} = 1 - 2\epsilon \ln(p) + \mathcal{O}(\epsilon^2) = 1 - \epsilon \ln(p^2) + \mathcal{O}(\epsilon^2)
\end{equation}
and
\begin{equation}
    \Gamma(\epsilon) = \frac{1}{\epsilon} - \gamma + \mathcal{O}(\epsilon)
\end{equation}
We get:
\begin{equation}
\begin{gathered}
    \Sigma_\Theta = -i\frac{\lambda^2}{3} \frac{I_{4 - \epsilon}}{2(4\pi)^{4 - \epsilon}} p^2 \frac{\frac{1}{\epsilon} - \gamma + \mathcal{O}(\epsilon)}{1 - \epsilon}\left( 1 - \epsilon \ln(p^2) + \mathcal{O}(\epsilon^2)\right) = \\ =
    -i\lambda^2 \frac{I_{4 - \epsilon}}{6(4\pi)^{4 - \epsilon}} p^2 \left(\frac{1}{\epsilon} - \gamma + \mathcal{O}(\epsilon)\right)\left( 1 - \epsilon \ln(p^2) + \mathcal{O}(\epsilon^2)\right)\left(1 + \epsilon\right) = \\ =
    -i\lambda^2 \frac{I_{4 - \epsilon}}{6(4\pi)^{4 - \epsilon}} p^2 \left(\frac{1}{\epsilon} + 1 - \ln(p^2) - \gamma + \mathcal{O}(\epsilon)\right) = \\ =
    -i\lambda^2 \frac{I_4}{6(4\pi)^4} p^2 \left(-\frac{1}{\epsilon} + \ln(\frac{p^2}{M^2}) + \mathcal{O}(\epsilon)\right)
\end{gathered}
\end{equation}
Where in the last step we took $\epsilon \rightarrow 0$ for the prefactor only and took the constants into the regularization constant $M^2$. Last thing left is to evaluate $I_4$.
\newline
This next section is a lot of algebra, which might not be very interesting, so if you don't care you can skip to the result which is half, at the end of the page 10.
\newline
First, let us write $I_4$:
\begin{equation}
    I_4 = \int_0^1 dx \int_0^{1-x} \frac{\left(\gamma - \gamma^2\right)dy}{(1 - y)^3(\beta + \alpha - \alpha^2)}
\end{equation}
Next, we simplify $\gamma - \gamma^2$ in terms of $\alpha$ and $\beta$:
\begin{equation}
    \gamma - \gamma^2 = \frac{\alpha - \alpha^2}{\beta + \alpha - \alpha^2} - \left(\frac{\alpha - \alpha^2}{\beta + \alpha - \alpha^2}\right)^2 = \frac{\beta\left(\alpha - \alpha^2\right)}{\left(\beta + \alpha - \alpha^2\right)^2}
\end{equation}
And plugging this back in, we get:
\begin{equation}
    I_4 = \int_0^1 dx \int_0^{1-x} \frac{\beta\left(\alpha - \alpha^2\right)dy}{(1 - y)^3(\beta + \alpha - \alpha^2)^3}
\end{equation}
Next, we write the fraction in terms of $x$ and $y$:
\begin{equation}
\begin{gathered}
    \frac{\beta\left(\alpha - \alpha^2\right)}{(1 - y)^3(\beta + \alpha + \alpha^2)^3} = 
    \frac{\frac{y}{1 - y}\left(\frac{1 - y - x}{1 - y} - \left(\frac{1 - y - x}{1 - y}\right)^2\right)}{(1 - y)^3\left(\frac{y}{1 - y} + \frac{1 - y - x}{1 - y} - \left(\frac{1 - y - x}{1 - y}\right)^2\right)^3} = \\ =
    \frac{\frac{y}{(1 - y)^3}\left((1 - y - x)(1 - y) - (1 - y - x)^2\right)}{\frac{(1 - y)^3}{(1 - y)^6}\left(y(1 - y) + (1 - y - x)(1 - y) - (1 - y - x)^2\right)^3} = \\ =
    \frac{xy(1 - y - x)}{\left((1 - x)(1 - y) - (1 - y - x)^2\right)^3}
\end{gathered}
\end{equation}
So, the integral becomes:
\begin{equation}
    I_4 = \int_0^1 dx \int_0^{1-x} \frac{xy(1 - y - x)dy}{\left((1 - x)(1 - y) - (1 - y - x)^2\right)^3}
\end{equation}
Now, it would be useful to write the $y$ integral from $0$ to $1$ instead of $1 - x$. So, let us define $y = (1 - x)u$ and $dy = (1 - x)du$ to get:
\begin{equation}
\begin{gathered}
    1 - y - x = 1 - x - (1 - x)u = (1 - x)(1 - u)
    \\
    (1 - x)(1 - y) = (1 - x)(1 - (1 - x)u)
\end{gathered}
\end{equation}
Thus:
\begin{equation}
\begin{gathered}
    I_4 = \int_0^1 dx \int_0^{1} \frac{x(1 - x)^3u(1 - u)du}{\left((1 - x)(1 - (1 - x)u) - (1 - x)^2(1 - u)^2\right)^3} = \\ =
    \int_0^1 dx \int_0^{1} \frac{x(1 - x)^3u(1 - u)du}{(1 - x)^3\left(1 - (1 - x)u - (1 - x)(1 - u)^2\right)^3} = \\ =
    \int_0^1 dx \int_0^{1} \frac{xu(1 - u)du}{\left(1 - (1 - x)u - (1 - x)(1 - u)^2\right)^3} = \\ =
    \int_0^1 dx \int_0^{1} du \frac{xu(1 - u)}{\left((1 - u + u^2)x + u - u^2\right)^3}
\end{gathered}
\end{equation}
Since the integation limits are no longer dependent on $x$, we can integrate over $x$ first:
\begin{equation}
    I_4 = \int_0^{1} \frac{u - u^2}{\left(1 - u + u^2\right)^3} du \int_0^1 \frac{x}{\left(x + \frac{u - u^2}{1 - u + u^2}\right)^3} dx
\end{equation}
Using the integral:
\begin{equation}
\begin{gathered}
    \int \frac{x}{(x + a)^3}dx = \{ z \equiv x + a \} = \int \frac{z - a}{z^3}dz = \int \frac{1}{z^2} - \frac{a}{z^3} dz = -\frac{1}{z} + \frac{a}{2z^2} + C = \\ =
    \frac{a - 2z}{2z^2} + C = -\frac{a + 2x}{2(a + x)^2} + C
\end{gathered}
\end{equation}
We can write:
\begin{equation}
\begin{gathered}
    I_4 = -\int_0^{1} \frac{u - u^2}{\left(1 - u + u^2\right)^3} du 
    \frac{\frac{u - u^2}{1 - u + u^2} + 2x}{2\left(\frac{u - u^2}{1 - u + u^2} + x\right)^2}\Bigg|_{x=0}^1 = \\ =
     -\int_0^{1} \frac{u - u^2}{\left(1 - u + u^2\right)^3} du 
    \left[
        \frac{\frac{u - u^2}{1 - u + u^2} + 2}{2\left(\frac{u - u^2}{1 - u + u^2} + 1\right)^2} - 
        \frac{\frac{u - u^2}{1 - u + u^2}}{2\left(\frac{u - u^2}{1 - u + u^2}\right)^2}
    \right] = \\ =
     -\frac{1}{2}\int_0^{1} \frac{u - u^2}{\left(1 - u + u^2\right)^3} du 
    \left[
        (1 - u + u^2)\frac{(u - u^2) + 2(1 - u + u^2)}{\left((u - u^2) + (1 - u + u^2)\right)^2} - \frac{1 - u + u^2}{u - u^2}
    \right] = \\ =
     -\frac{1}{2}\int_0^{1} \frac{u - u^2}{\left(1 - u + u^2\right)^2} du 
    \left[
        (2 - u + u^2) - \frac{1}{u - u^2}
    \right] = \\ =
    -\frac{1}{2} \int_0^{1} du \left[
        \frac{u - u^2}{\left(1 - u + u^2\right)^2} +
        \frac{u - u^2}{1 - u + u^2} -
        \frac{1}{\left(1 - u + u^2\right)^2}
    \right] = \\ =
    -\frac{1}{2} \int_0^{1} du \left[
        \frac{u - u^2 - 1}{\left(1 - u + u^2\right)^2} +
        \frac{u - u^2}{1 - u + u^2}
    \right] = \\ =
    -\frac{1}{2} \int_0^{1} du \left[
        \frac{-1}{1 - u + u^2} +
        \frac{u - u^2}{1 - u + u^2}
    \right] = \frac{1}{2}\int_0^1 du = \frac{1}{2}
\end{gathered}
\end{equation}
Plugging that into our result we get, as expected:
\begin{equation}
    \Sigma_\Theta = -i\lambda^2 \frac{1}{12(4\pi)^4} p^2 \left(-\frac{1}{\epsilon} + \ln(\frac{p^2}{M^2}) + \mathcal{O}(\epsilon)\right)
\end{equation}
\subsection{Taking $p^2 \rightarrow 0$}
Generally, for a particle of mass $m$ there is a pole at $p^2 = m^2$. In this case, for $m^2 = 0$, we expect a pole at $p^2 = 0$. While the zeroth order propegator has a pole at that point, this diagram contributes, not another pole, but a branch cut.
\end{document}
